\documentclass[12pt]{article}
\usepackage[utf8]{inputenc}
\usepackage[english]{babel}

\usepackage[letterpaper, margin = 1in]{geometry}
\usepackage{fancyhdr}
\usepackage{titling}
\usepackage{titlesec}
\usepackage{hyperref}
\usepackage{lastpage}
\usepackage{lipsum}

\usepackage[nocenter]{qtree}
\usepackage{tree-dvips}
\usepackage{linguex}

\usepackage[backend=biber, style=apa]{biblatex}
\addbibresource{deverbals.bib}
\usepackage{csquotes}
\DeclareLanguageMapping{english}{english-apa}

\usepackage{amssymb}
\usepackage{hanging}
\usepackage{tabularx}
\usepackage{multirow}
\usepackage{array}
\usepackage{multirow}
\renewcommand\tabularxcolumn[1]{m{#1}}

\usepackage{gb4e}

\title{CSCL 1301W: Syllabus and Assignment Materials}
\author{}
\date{}

\makeatletter
\renewcommand{\maketitle}
{
\begin{flushleft}
\textbf{{\large\@title}}
\newline
\@author
\end{flushleft}
}
\makeatother

\pagestyle{fancy}
\fancyhf{}
\lhead{Ian Thomas White}
\chead{\small{CSCL1301W, Spring 2020}}
\rhead{\thepage \ of \pageref{LastPage}}

\widowpenalty=100000
\clubpenalty=100000

\titleformat*{\section}{\medium\bfseries}
\titleformat*{\subsection}{\small\bfseries}
\titleformat*{\subsubsection}{\small\bfseries}

\makeatletter
\renewcommand\part{%
  \par\bigskip
  \secdef\@part\@spart}
\def\@part[#1]#2{%
  \refstepcounter{part}%
  \addcontentsline{toc}{part}{#1}%
  {\noindent\large\bfseries #2\par}%
  \nobreak\smallskip}
\def\@spart#1{%
  {\noindent\large\bfseries #1\par}%
  \nobreak\smallskip}
\makeatother

\begin{document}

\thispagestyle{empty}

\begin{flushright}
\small{CSCL1301W: \textbf{Reading Culture, Spring 2020} \\ instructor: \textbf{Ian White} \\ class setting: \textbf{MWF 11:15-12:05, Nicholson 315} \\ office hours: \textbf{by appointment} \\ contact: \textbf{whit2074@umn.edu}}
\end{flushright}

\maketitle

\tableofcontents 
\newpage

\renewcommand{\labelitemi}{$\square$}

\part{Syllabus}

\section{Course Description}
Rami Malek famously said, on \textit{The Tonight Show}, ``What the fuck did language ever do for me?''. This course will try and fail to figure out what language ever did for Rami Malik, for the remaining cast of \textit{Night at the Museum}, and for the full host of our dear \textit{homo sapiens}. That is, we will press upon our temples the unruly question of language's societal function, and we will sleuth out what one might glean about a particular society form its texts. Therein, we will examine four major moments of high-order (i.e., not strictly technical) linguistic theory, two (relatively) ancient, two (relatively) modern. Still, linguistic theory has its limits. To more comprehensively reckon our language's role, we will also pit against each linguistic theory a set of lyric objects from the same respective time period. These objects comprise poetry, music, and film, but all consider relatively directly the character and effects of language, including their own. As a more general and abstract goal, by analyzing and discussing these texts, we will work on how to form and argue precise arguments in speech and writing. In fact, this course meets the following LE requirements: Arts/Humanities, Writing Intensive.

\section{Texts}
We will be reading the following texts, all of which will be required, per the schedule. All texts are available in .pdf format on Canvas; although, of course, you're welcome to print or purchase hard copies (available at the Coffman bookstore).

\subsection{Linguistic Theory}
\begin{hangparas}{.25in}{1} \small{
    Chomsky, Noam. \textit{Cartesian Linguistics}. Cambridge, MA: Cambridge UP, 2009.
    
    ---. \textit{Syntactic Structures}. Berlin: Mouton de Gruyter, 2002.
    
    Plato. \textit{Cratylus}, in \textit{Complete Works}. Trans. by C.D.C. Reeve. Ed. by John Cooper. Indianapolis, IN: Hackett, 1997.
    
    de Saussure, Ferdinand. \textit{Course in General Linguistics}. Trans. by Roy Harris. London: Bloomsbury, 2013.
    
    Sextus Empiricus. ``Against the Grammarians'', in \textit{Against the Disciplines}. Trans. by Richard Bett. Oxford: Oxford UP, 2018.
    
}
    
\end{hangparas}
\noindent

\subsection{Lyric}

\begin{hangparas}{.25in}{1} \small{
    Brakhage, Stan. \textit{Cat's Cradle}. 1959. [film]
    
    Callimachus. ``Hymn to Apollo'', ``Hymn to Artemis'', in \textit{The Hymns}. Trans. by Susan Stephens. Oxford: Oxford UP, 2015.
    
    Debussy, Claude. \textit{Pr\'{e}lude \`{a} l'apr\`{e}s-midi d'un faune}. [audio]
    
    Deren, Maya. \textit{Meshes of the Afternoon}. 1943. [film]
    
    Dickinson, Emily. 353, 516, 543, 544, 883, 943, 1084, 1115, 1129, in \textit{The Complete Poems}. Ed. by Thomas Johnson. Boston: Little, Brown and Company, 1960.
    
    Lorde, Audre. ``Coal'', in \textit{The Collected Poems}. New York: W. W. Norton and Company, 1997.
    
    Mallarm\'{e}, St\'{e}phane. ``The Afternoon of a Faun'', in \textit{Collected Poems}. Trans. by Henry Weinfeld. Berkeley: University of California Press, 2011. 
    
    Marker, Chris. \textit{La Jet\'{e}e}. 1962. [film]
    
    Pindar. ``Olympian 1'', ``Olympian 7'', in \textit{Odes}. Trans. by Andrew Miller. Berkeley: University of California Press, 2019. 
    
    Plath, Sylvia. ``Daddy'', in \textit{ Collected Poems}. New York: HarperCollins Publishers, 1992.
    
    Stein, Gertrude. ``If I Told Him, A Complete Portrait of Picasso'', in \textit{Selections: Gertrude Stein}. Berkeley: University of California Press, 2008.
  
    ---. ``Stanza 1'', in \textit{Stanzas in Meditation and Other Poems}. Los Angeles: Sun and Moon, 1994.
    
    ---. ``Susie Asado'', in \textit{Selected Writings of Gertrude Stein}. New York: Peter Smith, 1992.
    
    Sun Ra. \textit{Antique Blacks}. 1974. [audio]
}
\end{hangparas}
\noindent

\section{Requirements}

\subsection{Reading and Notes}
The bulwark work of this class is to read our objects and to read them well. Reading, however, is a notoriously difficult task w/r/t both metabolism and second-order rigor. To add insult to injury, the objects we will read include various pitfalls and rough patches. To facilitate the understanding of, and the capacity to understand, complex texts, we will read our objects in short spurts, with the intention of being precise in our analysis therein. Most weeks have about 30 pages of reading, with Mondays and Wednesdays dedicated to the linguistic theory texts, Fridays to the lyric texts. Moreover, every reading assignment requires accompanying textual notes. Namely, you will need to upload to Canvas, before class start, pictures or screenshots of your textual apparatus marked-up, including, e.g.: questions you have while reading, definitions of unfamiliar words, underlinings of important portions, connections to other ideas, etc. Or, you could upload notes in a separate document, organized chronologically by page/section. These tasks should, even outside of this class, accompany any instance of reading, so the goal here is to develop that habit.

\subsection{Participation}
The objects of this class, as mentioned, are difficult, an obstacle significantly mollified by the input of others. Participation, then, is also required, that we might sap as much as possible from our objects. Note, as well, that overly placid classes are awkward for everyone involved. For those for whom speaking in class is an anxiety, feel free to meet with me separately to discuss the texts, instead or in addition. Participation points can potentially be made up with short writing assignments, upon my discretion.

\subsection{Outlines and Claims}
In addition to the textual notes, there are two separate, small but consistent assignments. First, each portion of linguistic theory reading requires an outline of the text's argument, an outline which details the main arguments, sub-arguments, assumptions, and evidences being made or invoked. These will be due on Wednesdays, comprising the Monday and Wednesday readings. Group work is allowed for these, but each person should upload their own copy that notes all group members. Second, each person must pick one small portion (1-2 sentences) from each lyric object and make an arguable claim about that portion w/r/t the object as a whole. These will be due on Fridays and should be done individually.

\subsection{Close Reading Papers}
Two short, concise papers are required, each two double-spaced pages, due on Canvas by 5 p.m. on, respectively 03/06 and 04/10. Each is based on one of the weekly claims made about a lyric object. The goal is to develop a rigorous and sound argument based on a hyper-close look at a small amount of text. Both of these papers can be revised for a better grade, due on 05/08. More information on how to write a close reading paper will be provided throughout the course.

\subsection{Final Paper} 
The final paper is longer, 5-6 double-spaced pages, due on Canvas by 5 p.m. on 05/08. This paper will pit any one linguistic text against any one lyric text, in order to discuss how they enhance and/or contradict each other. Use of one of the weekly lyric claims is fair game, here, as well. The goal is to route the type of close, material analysis utilized in the first paper into a sustained analysis of two texts together. This paper cannot be revised for a better grade. More information on how to conduct a sustained textual argument will be provided throughout the course.

\section{Grades}
There are 500 total points available, and your grade will be calculated as X/500, according to the following break down:
%\renewcommand{\labelitemi}{$\cdot$}
\begin{itemize}
    \item Reading Notes = 100 points. There are 33 days for which reading notes will be due, and I'll give you 3 ``free'' days for which you can decide not to upload notes, without penalty. Therefore, there are 30 gradable days each worth 100/30 ($\approx$3.33) points. The grading, though, will essentially be binary, i.e., either 100/30 or 0 points.
    \item Participation = 100 points. There are 37 days for which participation will be graded, and I'll give you 5 ``free'' days for which you can decide not to participate, without penalty (including not coming to class). Therefore, there are 32 gradable days each worth 100/32 (=3.125) points. The grading, as for the notes, will essentially be binary, i.e., either 100/32 or 0 points.
    \item Outlines/Claims = 100 points. There are 19 days for which an outline or a claim will be due, and I'll give you 1 ``free'' day, for which you can decide not to upload either, without penalty. Therefore, there are 18 gradable days, each worth 100/18 ($\approx$5.56) points.
    \item Close Reading Papers = 100 points. Each paper is worth 50 points. These papers will be graded with strict attention to argument and language and will then be given a square-root curve. The grade of a revised paper will replace that of the original.
    \item Final Paper = 100 points. This paper will be graded with strict attention to argument and language and will then be given a square-root curve.
\end{itemize}

\section{Policies and Resources}
\subsection{University-Wide Policies}
See the following links for relevant policies: \href{http://regents.umn.edu/sites/regents.umn.edu/files/policies/Student_Conduct_Code.pdf}{\underline{Student Conduct}}, \href{https://regents.umn.edu/sites/regents.umn.edu/files/policies/Sexual_Harassment_Sexual_Assault_Stalking_Relationship_Violence.pdf}{\underline{Sexual Harassment}}, \href{http://regents.umn.edu/sites/regents.umn.edu/files/policies/Equity_Diversity_EO_AA.pdf}{\underline{Equity+Diversity}}, \href{https://communitystandards.umn.edu/avoid-violations/avoiding-scholastic-dishonesty}{\underline{Scholastic Dishonesty}}, \href{https://regents.umn.edu/sites/regents.umn.edu/files/policies/Academic_Freedom.pdf}{\underline{Freedom+Responsibility}}

\subsection{Course-Specific Policies}
\begin{itemize}
    \item Any scholastic dishonesty will result in a failing grade on the assignment in question.  Don’t plagiarize. I often spot it, and it puts everyone involved in an annoying and pointless ethical situation. If you’re having trouble with a paper or other assignment, please don't hesitate to speak with me.
    \item A hostile environment will not be tolerated. Hostile behavior will result in the offending student being asked to leave for the day and to consult with me before being allowed to attend subsequent classes. All days so missed will result in a 0 for participation/attendance.
    \item Feel free to speak with me about any problems you're having with the course---we can work at least most things out. Note that I am required as a university employee to report certain cases related to Title IX violations and to mental health issues.
    \item If you are or will be registered with the DRC, please let me know as soon as possible so we can figure out how your accommodations will work for this course.
\end{itemize}

\subsection{University-Wide Resources}
\begin{itemize}
    \item The Center for Writing: \href{http://writing.umn.edu/}{writing.umn.edu}, 612-626-7579, 10 Nicholson Hall
    \item International Student Services: \href{https://isss.umn.edu/}{isss.umn.edu}, 612-626-7100, 190 Hubert Humphrey Center
    \item University Counseling Services: \href{https://counseling.umn.edu/}{counseling.umn.edu}, 612-624-3323, 340 Appleby Hall

    \item Disability Services: \href{https://disability.umn.edu/}{disability.umn.edu}, 612-626-1333, 80 McNamara Alumni Center
    \item Equal Opportunity / Affirmative Action: \href{https://eoaa.umn.edu/}{eoaa.umn.edu}, eoaa@umn.edu, 274 McNamara Alumni Center, 612-624-9547
    \item Boynton Health Services: \href{https://boynton.umn.edu/}{boynton.umn.edu}, 612-625-8400, 410 Church Street
    \item Boynton Mental Health Clinic: \href{http://mentalhealth.umn.edu/}{mentalhealth.umn.edu}, 612-624-1444, 410 Church Street
    \item Aurora Center: \href{http://aurora.umn.edu/}{aurora.umn.edu}, 612-626-2929, 24-hr hotline: 612-626-9111, 117 Appleby Hall
    \item Campus Police: \href{https://publicsafety.umn.edu/}{publicsafety.umn.edu/}, non-emergencies: 612-624-2677, emergencies: 911, Transportation Safety Building (Washington Avenue)
\end{itemize}

\newpage

\section{Schedule}

\begin{center}
\begin{tabularx}{\textwidth}{| c 
 | >{\centering\arraybackslash}X 
 | >{\centering\arraybackslash}X
 | >{\centering\arraybackslash}X |}
    \hline
    \footnotesize{Week} & \footnotesize{Monday} & \footnotesize{Wednesday} & \footnotesize{Friday} \\
    \hline \hline
    \footnotesize{[all] text assignments} & \footnotesize{due: theory textual notes} & \footnotesize{due: theory textual notes + theory outlines} & \footnotesize{due: lyric textual notes + lyric claims} \\
    \hline \hline
    \footnotesize{[1] 01/20-01/24} & \footnotesize{no class} & \footnotesize{intro} & \footnotesize{grammar/writing review} \\
    \hline
    \footnotesize{[2] 01/27-01/31} & \footnotesize{grammar/writing review [worksheet required]} & \footnotesize{Greek language basics; Herodotus / Shakespeare [notes not required]} & \footnotesize{Greek background; \textit{Cratylus}, 383a-385e} \\
    \hline \hline
    \footnotesize{[3] 02/03-02/07} & \footnotesize{\textit{Cratylus}, 385e-396e} & \footnotesize{\textit{Cratylus}, 396e-420e} & \footnotesize{Pindar, ``Olympian 1''} \\
    \hline
    \footnotesize{[4] 02/10-02/14} & \footnotesize{\textit{Cratylus}, 420e-435d} & \footnotesize{\textit{Cratylus}, 435d-440e} & \footnotesize{Pindar, ``Olympian 7''} \\
    \hline \hline
    \footnotesize{[5] 02/17-02/21} & \footnotesize{``Against the Grammarians'', [1]-[40]} & \footnotesize{``Against the Grammarians'', [41]-[96]} & \footnotesize{Callimachus, ``Hymn to Apollo''} \\
    \hline
    \footnotesize{[6] 02/24-02/28} & \footnotesize{``Against the Grammarians'', [97]-[175]} & \footnotesize{``Against the Grammarians'', [176]-[247]} & \footnotesize{Callimachus, ``Hymn to Artemis''} \\
    \hline
    \footnotesize{[7] 03/02-03/06} & \footnotesize{``Against the Grammarians'', [296]-[320]} & \footnotesize{Saussure background [outline optional]} & \footnotesize{no class; \textbf{due: Close Reading 1}} \\
    \hline\hline
    \footnotesize{[8] 03/09-03/13} & \footnotesize{no class} & \footnotesize{no class} & \footnotesize{no class} \\
    \hline \hline
    \footnotesize{[9] 03/16-03/20} & \footnotesize{\textit{Course in General Linguistics}, 1-36 [notes/outline optional], 73-116} & \footnotesize{\textit{Course in General Linguistics}, 117-130; Debussy, ``Prelude to Afternoon of a Faun''} & \footnotesize{Mallarm\'{e}, ``The Afternoon of a Faun'; Debussy, ``Prelude to Afternoon of a Faun''} \\
    \hline
    \footnotesize{[10] 03/23-03/27} & \footnotesize{\textit{Course in General Linguistics}, 131-148} & \footnotesize{\textit{Course in General Linguistics}, 149-164; Dickinson, 1085} & \footnotesize{Dickinson, 353, 516, 543, 544, 883, 943, 1084, 1115, 1129} \\
    \hline
    \footnotesize{[11] 03/30-04/03} & \footnotesize{\textit{Course in General Linguistics}, 165-198} & \footnotesize{\textit{Course in General Linguistics}, 199-224; Stein, ``If I Told Him''} & \footnotesize{Stein, ``Stanza 1'', ``Susie Asado'', ``If I Told Him''} \\
    \hline \hline
    \footnotesize{[12] 04/06-04/10} & \footnotesize{\textit{Cartesian Linguistics}, 55-77} & \footnotesize{\textit{Cartesian Linguistics}, 78-108} & \footnotesize{Chomsky background; \textbf{due: Close Reading 2}} \\
    \hline
    \footnotesize{[13] 04/13-04/17} & \footnotesize{\textit{Syntactic Structures}, 5-25} & \footnotesize{\textit{Syntactic Structures}, 26-48; Brakhage, \textit{Cat's Cradle}} & \footnotesize{Deren, \textit{Meshes of the Afternoon}; Brakhage, \textit{Cat's Cradle}} \\
    \hline
    \footnotesize{[14] 04/20-04/24} & \footnotesize{\textit{Syntactic Structures}, 49-60} & \footnotesize{\textit{Syntactic Structures}, 61-84} & \footnotesize{Plath, ``Daddy''; Lorde, ``Coal''} \\
    \hline
    \footnotesize{[15] 04/27-05/01} & \footnotesize{\textit{Syntactic Structures}, 85-108} & \footnotesize{Marker, \textit{La Jet\'{e}e}; Sun Ra, \textit{Antique Blacks} [outline optional]} & \footnotesize{wrap up} \\
    \hline \hline
    \footnotesize{[16] 05/04-05/08} & \footnotesize{no class} & \footnotesize{no class} & \footnotesize{no class; \textbf{due: Final Paper; all revisions}} \\
    \hline
\end{tabularx}
\end{center}

\newpage
\part{Assignment Guidelines and Rubrics}

\section{Theory Outlines}

\subsection{Guidelines}

The point of the outlines is to figure out and present the arguments/claims of each relevant text in a relatively concise manner. That is, the outline should consist, in normal outline fashion, of a set of ranked claims such that: the primary claim(s) are positioned at the outer most level and all supporting claims/proofs/examples at progressively inward levels. Outlines should include reference numbers comprising the relevant start and endpoint of any given claim (for film/audio, use timestamps), which means that inner argument references are subsets of their controlling outer arguments. Good outlines should include, if relevant, the types of arguments being made (e.g. `proof by analogy', `assumption', `example', etc.).

\renewcommand{\labelenumi}{(\Roman{enumi})}
\renewcommand{\labelenumii}{(\Alph{enumii})}
\renewcommand{\labelenumiii}{(\roman{enumiii})}
\renewcommand{\labelenumiV}{(\alph{enumiV})}
\renewcommand{\labelenumV}{(\alph{enumV})}

\subsection{Format}
\begin{small}
\begin{enumerate}
\item \textbf{first main argument}, i.e., a primary (non-subordinated) claim to be shown, per the given section (could be carried over from a prior section); \textbf{location X-Y}, e.g., pp. 1-20
    \begin{enumerate}
    \item \textbf{first sub-argument}, i.e., some supporting evidence for the main argument (e.g., examples, smaller argument to be proven, etc.); \textbf{location [X+a]-[Y--b]}, e.g., pp. 3-12
        \begin{enumerate}
        \item \textbf{first sub-sub-argument}, i.e., some supporting evidence for the sub-argument; \textbf{location [[X+a]+c]-[[Y--b]--d]}, e.g., pp. 3-5
        \item \textbf{second sub-sub-argument}, i.e., some new/different supporting evidence for the sub-claim; \textbf{location [[X+a]+e]-[[Y--b]--f]}, e.g., pp. 5-9
        \item ... [and so on, including further subordination, until the end of the first sub-argument, i.e., location Y--b, here, p. 12]
        \end{enumerate}
    \item \textbf{second sub-argument}, i.e., some new/different supporting evidence for the main argument; \textbf{location [X+g]-[Y--h]}, e.g., pp. 13-17
    \item ... [and so on until the end of the first main argument, i.e., location Y, here, p. 20]
    \end{enumerate}
\item \textbf{second main argument}, i.e., a new/different non-subordinated argument; \textbf{location W-Z} (e.g., pp. 21-35)
\item ... [and so on until the end of the reading section]
\end{enumerate}
\end{small}

\subsection{Rubric}
\begin{center}
\begin{small}
\begin{tabularx}{\textwidth}{
| >{\centering\arraybackslash}X
| >{\centering\arraybackslash}X
|| >{\centering\arraybackslash}X 
| >{\centering\arraybackslash}X
| >{\centering\arraybackslash}X
|| >{\centering\arraybackslash}X
|}
\hline
completion & format & accuracy & completeness & precision & \textbf{grade} \\
\hline
assignment was turned in on time & arguments are ranked with textual locations & arguments and descriptions are correct w/r/t the text & all arguments from reading section are included & arguments are minimally divided and described well  & sum of preceding sections \\
\hline
\multicolumn{1}{|r|}{/ 1.56} & \multicolumn{1}{r||}{/ 1} & \multicolumn{1}{r|}{/ 1} & \multicolumn{1}{r|}{/ 1} & \multicolumn{1}{r||}{/ 1} & \multicolumn{1}{r|}{/ 5.56} \\
\hline
\multicolumn{5}{c|}{} & \multicolumn{1}{c|}{=} \\
\cline{6-6}
\end{tabularx}
\end{small}
\end{center}

\newpage

\section{Lyric Claims}

\subsection{Guidelines}

One of the most difficult parts of writing papers and making arguments is coming up with a claim about a text. A claim is not a summary or a simple observation of a text. Rather, a claim is a succinct statement about what a text, or some aspect of a text, is doing, to be proven using once certain textual features. The claims due on Friday will be used for the first two close reading papers. We will go over some of these claims in class, to practice. See the grammar review handout for more information on thesis claims.

\subsection{Rubric}

\begin{center}
\begin{small}
\begin{tabularx}{\textwidth}{
| >{\centering\arraybackslash}X
| >{\centering\arraybackslash}X
|| >{\centering\arraybackslash}X 
| >{\centering\arraybackslash}X
| >{\centering\arraybackslash}X
|| >{\centering\arraybackslash}X
|}
\hline
completion & format & accuracy & completeness & precision & \textbf{grade} \\
\hline
assignment was turned in on time & claim is more than just summary & claim makes sense w/r/t the text & claim is a fully recognizable proposition & claim encapsulates well the portion of text  & sum of preceding sections \\
\hline
\multicolumn{1}{|r|}{/ 1.56} & \multicolumn{1}{r||}{/ 1} & \multicolumn{1}{r|}{/ 1} & \multicolumn{1}{r|}{/ 1} & \multicolumn{1}{r||}{/ 1} & \multicolumn{1}{r|}{/ 5.56} \\
\hline
\multicolumn{5}{c|}{} & \multicolumn{1}{c|}{=} \\
\cline{6-6}
\end{tabularx}
\end{small}
\end{center}
\newpage

\section{Close Reading Papers}

\renewcommand{\labelitemi}{$\square$}

\subsection{Guidelines}

Your first two papers are argumentative essays, each based on one of your Friday lyric claims. That is, the claim acts as the thesis that the paper seeks to prove, using precise textual evidence. If needed, you may modify the original claims to fit the papers, but the scope of the claim should remain quite small (i.e. a few lines/sentences). The point here is to be precise and rigorous, rather than expansive. You may revise each paper once, but revisions are not required. \textbf{Due: All by 5.pm. on Canvas; Paper 1: 03/06; Paper 2: 04/10; revisions: 05/08}

\subsection{Format}

\begin{itemize}
\item 2 pages, double-spaced, 12 pt. TNR font, 1 in. margins, titled descriptively, \textbf{.pdf only}
\item include in-text citations and works cited (any format), 0 outside sources allowed
\end{itemize}

\subsection{Rubric, First Draft}

\begin{center}
\begin{small}
\begin{tabularx}{\textwidth}{
| >{\centering\arraybackslash}X
|| >{\centering\arraybackslash}X
| >{\centering\arraybackslash}X 
| >{\centering\arraybackslash}X
| >{\centering\arraybackslash}X
|| >{\centering\arraybackslash}X
|}
\hline
format & accuracy & completeness & insight & precision & \textbf{grade} \\
\hline
paper has correct format and length and is checked for grammar and spelling & arguments and descriptions are correct and relevant w/r/t the text & all arguments made that are necessary for the claim are included and none are extraneous & claim and arguments aren't superficial or summary of text or class discussion & arguments are clear, well-formed, and logically minimal (i.e., no skipped steps) & square root of the sum of preceding sections reconverted to points out of 50 \\
\hline
\multicolumn{1}{|r||}{/ 10} & \multicolumn{1}{r|}{/ 10} & \multicolumn{1}{r|}{/ 10} & \multicolumn{1}{r|}{/ 10} & \multicolumn{1}{r||}{/ 10} & \multicolumn{1}{r|}{$\sqrt{\hspace{2em} / 50}$} \\
\hline
\multicolumn{5}{c|}{} & \multicolumn{1}{c|}{=} \\
\cline{6-6}
\multicolumn{5}{c|}{} & \multicolumn{1}{c|}{$\rightarrow$ \hfill / 50} \\
\cline{6-6}
\multicolumn{5}{c|}{} & \multicolumn{1}{c|}{=} \\
\cline{6-6}
\end{tabularx}
\end{small}
\end{center}

\subsection{Rubric, Second Draft (Optional)}
\begin{center}
\begin{small}
\begin{tabularx}{\textwidth}{
| >{\centering\arraybackslash}X
|| >{\centering\arraybackslash}X
| >{\centering\arraybackslash}X 
| >{\centering\arraybackslash}X
| >{\centering\arraybackslash}X
|| >{\centering\arraybackslash}X
|}
\hline
format & accuracy & completeness & insight & precision & \textbf{grade} \\
\hline
paper has correct format and length and is checked for grammar and spelling & arguments and descriptions are correct and relevant w/r/t the text & all arguments made that are necessary for the claim are included and none are extraneous & claim and arguments aren't superficial or summary of text or class discussion & arguments are clear, well-formed, and logically minimal (i.e., no skipped steps) & square root of the sum of preceding sections reconverted to points out of 50 \\
\hline
\multicolumn{1}{|r||}{/ 10} & \multicolumn{1}{r|}{/ 10} & \multicolumn{1}{r|}{/ 10} & \multicolumn{1}{r|}{/ 10} & \multicolumn{1}{r||}{/ 10} & \multicolumn{1}{r|}{$\sqrt{\hspace{2em} / 50}$} \\
\hline
\multicolumn{5}{c|}{} & \multicolumn{1}{c|}{=} \\
\cline{6-6}
\multicolumn{5}{c|}{} & \multicolumn{1}{c|}{$\rightarrow$ \hfill / 50} \\
\cline{6-6}
\multicolumn{5}{c|}{} & \multicolumn{1}{c|}{=} \\
\cline{6-6}
\end{tabularx}
\end{small}
\end{center}

\newpage

\section{Final Paper, New Option}

\renewcommand{\labelitemi}{$\square$}

\subsection{Guidelines}

Like the two close reading papers, this final paper is an argumentative paper, based on one of the linguistic theory texts. The paper should have a thesis claim that the paper seeks to prove, using precise textual evidence. The task here is to pick apart at most one paragraph from the one linguistic theory text. Some potential questions to ask: How does language function? What is language's societal role? Is language beneficial/harmful, simple/complex, understood/misunderstood? This shouldn't just be a summary of the paragraph or the text itself but should rather draw precise implications from the text. As in the other papers, the point here is to be precise and rigorous, rather than expansive. \textbf{Due: All by 5.pm. on Canvas; Paper 1: 03/06; Paper 2: 04/10; revisions: 05/08}

\subsection{Format}

\begin{itemize}
\item 3 pages, double-spaced, 12 pt. TNR font, 1 in. margins, titled descriptively, \textbf{.pdf only}
\item include in-text citations and works cited (any format), 0 outside sources allowed
\end{itemize}

\subsection{Rubric, First Draft}

\begin{center}
\begin{small}
\begin{tabularx}{\textwidth}{
| >{\centering\arraybackslash}X
|| >{\centering\arraybackslash}X
| >{\centering\arraybackslash}X 
| >{\centering\arraybackslash}X
| >{\centering\arraybackslash}X
|| >{\centering\arraybackslash}X
|}
\hline
format & accuracy & completeness & insight & precision & \textbf{grade} \\
\hline
paper has correct format and length and is checked for grammar and spelling & arguments and descriptions are correct and relevant w/r/t the text & all arguments made that are necessary for the claim are included and none are extraneous & claim and arguments aren't superficial or summary of text or class discussion & arguments are clear, well-formed, and logically minimal (i.e., no skipped steps) & square root of the sum of preceding sections reconverted to points out of 50 \\
\hline
\multicolumn{1}{|r||}{/ 10} & \multicolumn{1}{r|}{/ 10} & \multicolumn{1}{r|}{/ 10} & \multicolumn{1}{r|}{/ 10} & \multicolumn{1}{r||}{/ 10} & \multicolumn{1}{r|}{$\sqrt{\hspace{2em} / 50}$} \\
\hline
\multicolumn{5}{c|}{} & \multicolumn{1}{c|}{=} \\
\cline{6-6}
\multicolumn{5}{c|}{} & \multicolumn{1}{c|}{$\rightarrow$ \hfill / 50} \\
\cline{6-6}
\multicolumn{5}{c|}{} & \multicolumn{1}{c|}{=} \\
\cline{6-6}
\end{tabularx}
\end{small}
\end{center}

\newpage

\section{Final Paper, Original Option}

\subsection{Guidelines}

Like the two close reading papers, the final paper is an argumentative paper, based on one lyric text and one linguistic theory text. Again, the paper should have a thesis claim that the paper seeks to prove, using precise textual evidence. In essence, the goal is to derive the implications of each text's explicit/implicit linguistic theories when those theories are taken together. Do the two theories contradict each other in a fundamental way? Does one include the other as a supposition or framework? Are both incomplete without the other? Note that your claim should be constructive; i.e., the paper is not just a compare-contrast essay. Rather, a precise argument should be made about the interactions of the two theories. You may use a modified/expanded version of one of your Friday lyric claims, but this is not required. In any case, the scope of the paper should still be relatively small, i.e., looking at a small amount of each of the two respective texts, and the paper should prioritize precision and rigor over expansiveness. Revision of this paper is not allowed. \textbf{Due: 05/08, 5 p.m., on Canvas}

\subsection{Format}

\begin{itemize}
\item 6 pages, double-spaced, 12 pt. TNR font, 1 in. margins, titled descriptively, \textbf{.pdf only}
\item include in-text citations and works cited (any format), 0 outside sources allowed
\end{itemize}

\subsection{Rubric}
\begin{center}
\begin{small}
\begin{tabularx}{\textwidth}{
| >{\centering\arraybackslash}X
|| >{\centering\arraybackslash}X
| >{\centering\arraybackslash}X 
| >{\centering\arraybackslash}X
| >{\centering\arraybackslash}X
|| >{\centering\arraybackslash}X
|}
\hline
format & accuracy & completeness & insight & precision & \textbf{grade} \\
\hline
paper has correct format and length and is checked for grammar and spelling & arguments and descriptions are correct and relevant w/r/t the text & all arguments made that are necessary for the claim are included and none are extraneous & claim and arguments aren't superficial or summary of text or class discussion & arguments are clear, well-formed, and logically minimal (i.e., no skipped steps) & square root of the sum of preceding sections reconverted to points out of 100 \\
\hline
\multicolumn{1}{|r||}{/ 20} & \multicolumn{1}{r|}{/ 20} & \multicolumn{1}{r|}{/ 20} & \multicolumn{1}{r|}{/ 20} & \multicolumn{1}{r||}{/ 20} & \multicolumn{1}{r|}{$\sqrt{\hspace{2em} / 100}$} \\
\hline
\multicolumn{5}{c|}{} & \multicolumn{1}{c|}{=} \\
\cline{6-6}
\multicolumn{5}{c|}{} & \multicolumn{1}{c|}{$\rightarrow$ \hfill / 100} \\
\cline{6-6}
\multicolumn{5}{c|}{} & \multicolumn{1}{c|}{=} \\
\cline{6-6}
\end{tabularx}
\end{small}
\end{center}

\end{document}
